\chapter{Accéder aux marchés~: taux et crédit}\label{ch:m2:getting-access-rates-credit}

Un swap se conclut en une seconde sur un écran. Avant le premier, un fonds a besoin d'un contrat-cadre
avec chaque dealer, négocié clause par clause, d'un \omterm{def:m2:swap-clearing-and-the-ccp-basis:umr}{credit support annex} qui précise quel collatéral
circule et quand, d'un identifiant qui indique aux régulateurs qui il est, et de comptes auprès d'un
\omterm{def:m2:getting-access-rates-credit:cb}{courtier compensateur} pour les swaps qui doivent être compensés. Le repo de Trésor exige son propre
contrat-cadre et, de plus en plus, un accès à la chambre de compensation~: les États-Unis imposent
la compensation du repo éligible à partir de juin 2027. Pour chaque élément, le fonds choisit entre
louer le bilan d'un tiers et devenir membre, ou construire le sien. Ce chapitre présente les documents,
le statut qu'une société a sur les plateformes de négociation, la manière dont les non-banques accèdent
aux marchés interdealers, le \omterm{def:m2:repo-and-specials:sponsored}{repo sponsorisé} et l'adhésion à une chambre de compensation, ainsi que les
\omterm{def:m2:getting-access-rates-credit:cb}{courtiers compensateurs} de swaps, et construit un modèle du coût de chaque voie.

\section{Documentation avant la première transaction}

\begin{definition}[Contrat-cadre ISDA, Global Master Repurchase Agreement]\label{def:m2:getting-access-rates-credit:isda}
Un \emph{contrat-cadre ISDA}\index{contrat-cadre ISDA} est le contrat type,
publié par l'International Swaps and Derivatives Association, qui régit l'ensemble des
dérivés de gré à gré que deux parties échangent entre elles~: un formulaire imprimé, une
annexe négociée entre elles et une confirmation pour chaque transaction, le tout formant
un seul accord. Le \emph{Global Master Repurchase Agreement}\index{Global Master Repurchase Agreement}
(GMRA) est la norme correspondante pour le repo, publiée par
l'International Capital Market Association.
\end{definition}

Le \omterm{def:m2:getting-access-rates-credit:isda}{contrat-cadre ISDA} existe dans les versions de 1992 et de 2002~; sa disposition centrale est
la compensation de clôture (close-out netting)~: en cas de défaut, chaque transaction en cours entre
les deux parties est remplacée par un montant unique de résiliation anticipée, dû par l'une à l'autre. C'est ce qui permet
au service crédit d'un dealer et à ses régulateurs de mesurer l'exposition nette plutôt que transaction par
transaction, à condition que la compensation soit opposable selon le droit de l'insolvabilité du pays de la
contrepartie. Le \omterm{def:m2:swap-clearing-and-the-ccp-basis:umr}{credit support annex} du chapitre 10 s'y superpose et
fixe le collatéral. Le GMRA, publié pour la première fois en 1992 et révisé pour la dernière fois en 2011, fait
de même pour le repo, avec des avis juridiques sur son opposabilité dans de nombreuses juridictions,
publiés par son promoteur (\cref{fig:m2:getting-access-rates-credit:docs}).

\begin{definition}[Identifiant d'entité juridique]\label{def:m2:getting-access-rates-credit:lei}
Un \emph{identifiant d'entité juridique}\index{identifiant d'entité juridique} (LEI) est un code unique
de 20 caractères, selon la norme ISO~17442, qui identifie une entité juridique dans
les déclarations de transactions et les dépôts réglementaires, avec des données de référence publiques sur ce qu'est l'entité
et sur qui la détient.
\end{definition}

Le LEI est né de la crise de 2008, lorsque les régulateurs ne pouvaient pas savoir qui devait quoi à
qui~: le G20 a approuvé la charte du système en 2012, et une fondation suisse, la GLEIF,
l'administre. Les régulateurs utilisent le LEI pour identifier les parties dans les déclarations de transactions, si bien qu'un
fonds en obtient un avant de négocier.

\begin{omfigure}
\centering
\begin{tikzpicture}[font=\small,
  doc/.style={draw,thick,rounded corners=2pt,align=center,minimum width=3.4cm,minimum height=0.75cm,inner sep=2pt}]
\node[doc,draw=omDef,fill=omDef!10] (ma) at (0,0) {contrat-cadre ISDA};
\node[doc,draw=omDef,fill=omDef!5] (sch) at (0,0.95) {annexe (négociée)};
\node[doc,draw=omThm,fill=omThm!8] (csa) at (0,1.9) {credit support annex};
\node[doc,draw=gray,fill=gray!8] (conf) at (0,2.85) {confirmations, par transaction};
\node[font=\footnotesize\bfseries] at (0,3.7) {swaps, par dealer};
\node[doc,draw=omProp,fill=omProp!10] (gm) at (5,0) {GMRA 2011};
\node[doc,draw=omProp,fill=omProp!5] (ann) at (5,0.95) {annexes (négociées)};
\node[doc,draw=gray,fill=gray!8] (rc) at (5,1.9) {confirmations, par transaction};
\node[font=\footnotesize\bfseries] at (5,3.7) {repo, par contrepartie};
\node[doc,draw=omThm,fill=omThm!8,minimum width=3.8cm] (clr) at (10,0) {accords de compensation};
\node[doc,draw=gray,fill=gray!8,minimum width=3.8cm] (plat) at (10,0.95) {onboarding sur les plateformes};
\node[doc,draw=omDef,fill=omDef!8,minimum width=3.8cm] (lei) at (10,1.9) {identifiant d'entité juridique};
\node[font=\footnotesize\bfseries] at (10,3.7) {une fois pour toutes};
\end{tikzpicture}

\omcaption{Les documents dont un fonds a besoin avant son premier swap et son premier repo~: un
\omterm{def:m2:getting-access-rates-credit:isda}{contrat-cadre ISDA} avec chaque dealer, accompagné de son annexe, du \omterm{def:m2:swap-clearing-and-the-ccp-basis:umr}{credit support annex} et des
confirmations~; un GMRA avec chaque contrepartie de repo~; et, une seule fois, les accords de compensation,
l'onboarding sur les plateformes et un LEI. Schéma.}\label{fig:m2:getting-access-rates-credit:docs}
\end{omfigure}

\section{Statut de dealer et de client sur les plateformes}

Les taux et le crédit se négocient sur deux types de plateformes électroniques (chapitres 4 et 22)~: les
marchés interdealers, où les dealers et quelques grandes non-banques négocient entre eux,
principalement sur des carnets d'ordres~; et les \omterm{def:m2:the-treasury-market:idb}{plateformes dealer-to-client}, où les clients demandent des cotations
aux dealers. Le statut d'une société détermine ce qu'elle voit et avec qui elle peut négocier~: en tant que
client, elle demande des cours~; en tant que dealer, elle répond aux demandes et peut voir des flux que les clients
ne voient pas.

Dans les swaps, les règles s'attaquent aux clubs fermés~: une \omterm{def:m2:interest-rate-swaps:sef}{swap execution facility} (SEF) américaine doit
donner à tout eligible contract participant un accès impartial à ses marchés, selon des critères
impartiaux, transparents et appliqués de manière équitable et non discriminatoire, avec des
commissions comparables pour un accès comparable. Avoir accès à une plateforme ne signifie pas encore
pouvoir négocier~: les contreparties présentes doivent accepter le crédit de la société, directement ou
par l'intermédiaire de son \omterm{def:m2:getting-access-rates-credit:cb}{courtier compensateur}.

\section{Accès aux marchés interdealers pour les non-banques}

Le marché interdealer du Trésor n'est plus un club de dealers. Le chapitre 4 a montré que, le
15 octobre 2014, les principal trading firms, des non-banques négociant pour leur propre compte,
représentaient plus de la moitié du volume des marchés interdealers au comptant et à terme
étudiés. Elles accèdent aux carnets d'ordres par l'onboarding propre aux plateformes et,
pour le crédit et le règlement-livraison, par l'intermédiaire d'une banque ou d'un membre compensateur qui se porte garant de leurs
transactions, selon le même schéma que le \omterm{def:m2:the-fx-market:sdp}{prime brokerage de change} du chapitre précédent. Pour une
non-banque, le carnet d'ordres donne un accès anonyme aux cours les plus serrés~; son prix
en est l'accord de crédit, ainsi que l'obligation de régler et de compenser ce qu'elle négocie.

\section{Repo sponsorisé et adhésion à une chambre de compensation}

\begin{definition}[Membre sponsorisé]\label{def:m2:getting-access-rates-credit:sponsored}
Un \emph{membre sponsorisé}\index{membre sponsorisé} est une société admise auprès d'une
chambre de compensation par l'intermédiaire d'un membre parrain, qui soumet les transactions du
membre sponsorisé à la compensation et garantit ses obligations envers la chambre de compensation.
\end{definition}

Le chapitre 5 a décrit le \omterm{def:m2:repo-and-specials:sponsored}{repo sponsorisé}. La Fixed Income Clearing Corporation, la chambre de compensation
du repo de Trésor, offre trois voies à un fonds
(\cref{fig:m2:getting-access-rates-credit:routes}). Il peut devenir \omterm{def:m2:getting-access-rates-credit:sponsored}{membre sponsorisé}
d'une banque qui est membre parrain~: en novembre 2024, la SEC a approuvé des modifications de règles
supprimant l'exigence que les \omterm{def:m2:getting-access-rates-credit:sponsored}{membres sponsorisés} soient des qualified institutional buyers, et les
banques parrains doivent disposer de fonds propres d'au moins 5~milliards de dollars et être
bien capitalisées. Il peut recourir au service de compensation par agent d'un membre, que la FICC présente désormais
comme un modèle «~done-away~»~: le fonds traite avec n'importe quelle contrepartie et l'agent compense. Ou il peut
demander à devenir membre lui-même, avec le capital, les opérations et les contributions au fonds de
compensation que cela suppose. Une quatrième voie, le \omterm{def:m2:repo-and-specials:triparty}{repo bilatéral} non compensé, se réduit à mesure que
l'\omterm{def:m2:swap-clearing-and-the-ccp-basis:mandate}{obligation de compensation} entre en vigueur.

\begin{omfigure}
\centering
\begin{tikzpicture}[font=\small,
  box/.style={draw,thick,rounded corners=2pt,align=center,minimum height=0.95cm,minimum width=2.8cm,inner sep=3pt},
  arr/.style={-{Stealth[length=2.5mm]},thick}]
\node[box,draw=omDef,fill=omDef!8] (f) at (0,0) {fonds};
\node[box,draw=omThm,fill=omThm!8] (s) at (5,1.6) {banque parrain};
\node[box,draw=omProp,fill=omProp!8] (a) at (5,0) {membre\\compensateur agent};
\node[box,draw=gray,fill=gray!15] (c) at (10.2,0) {FICC\\(chambre de compensation)};
\node[box,draw=gray,fill=gray!8] (d) at (5,-1.6) {dealer (bilatéral)};
\draw[arr] (f) -- node[above,sloped,pos=0.4,font=\footnotesize] {parrainage} (s);
\draw[arr] (s) -- (c);
\draw[arr] (f) -- node[above,font=\footnotesize] {done-away} (a);
\draw[arr] (a) -- (c);
\draw[arr,dashed] (f) -- node[below,sloped,pos=0.4,font=\footnotesize] {non compensé} (d);
\draw[arr,very thick,omDef] (f.north) to[out=60,in=150] node[pos=0.5,above=2pt,font=\footnotesize] {achat~: adhésion directe} (c.north);
\end{tikzpicture}

\omcaption{Voies d'accès au repo de Trésor compensé pour un fonds~: adhésion sponsorisée par une
banque, un membre compensateur agent qui compense les transactions réalisées avec des tiers, adhésion
directe à la chambre de compensation, et \omterm{def:m2:repo-and-specials:triparty}{repo bilatéral} non compensé avec un dealer, que l'\omterm{def:m2:swap-clearing-and-the-ccp-basis:mandate}{obligation de
compensation} réduit. Schéma indicatif.}\label{fig:m2:getting-access-rates-credit:routes}
\end{omfigure}

\begin{dated}{2026-09}{Compensation du Trésor et accès sponsorisé}\label{dat:m2:getting-access-rates-credit:access}
Règle de compensation du Trésor de la SEC~: dates de mise en conformité le 31 décembre 2026 pour les
transactions au comptant éligibles et le 30 juin 2027 pour les repos éligibles (chapitre~4). Le Sponsored Service de la FICC~:
volume quotidien moyen de 2.3~billions de dollars au deuxième trimestre 2026, selon la DTCC. Approbation par la
SEC des modifications des règles d'accès de la FICC, dont la suppression de l'exigence
de qualified institutional buyer pour les \omterm{def:m2:getting-access-rates-credit:sponsored}{membres sponsorisés}~: 21 novembre 2024.
\end{dated}

\section{Courtiers compensateurs de swaps}

\begin{definition}[Courtier compensateur (swaps)]\label{def:m2:getting-access-rates-credit:cb}
Un \emph{courtier compensateur}\index{courtier compensateur (swaps)} (clearing broker) pour les swaps est un membre compensateur
d'une chambre de compensation de swaps qui compense les swaps de ses clients, verse leur marge à la
chambre et en garantit l'exécution envers elle~; aux États-Unis, c'est un
futures commission merchant (FCM) enregistré pour les swaps.
\end{definition}

Les swaps soumis à l'\omterm{def:m2:swap-clearing-and-the-ccp-basis:mandate}{obligation de compensation} du chapitre~10 doivent être compensés~; un fonds a donc besoin
d'un \omterm{def:m2:getting-access-rates-credit:cb}{courtier compensateur} auprès des chambres de compensation où ses dealers souhaitent traiter avec lui. Le courtier
facture des commissions sur le notionnel ou par ticket, exige une marge initiale, souvent assortie
d'un supplément au titre du risque du fonds, et limite les positions du fonds selon son propre
appétit de crédit. Devenir directement membre compensateur est possible pour les grandes sociétés, mais
entraîne des contributions au fonds de défaillance, des obligations de gestion des défaillances et des
exigences opérationnelles que seuls les gros portefeuilles justifient.

\begin{proposition}[Quand acquérir l'accès]\label{prop:m2:getting-access-rates-credit:breakeven}
Soit une voie dont le coût annuel fixe est $F$ et le coût annuel variable $u$ par unité de
solde de repo, pour un portefeuille dont le notionnel de swaps est un multiple fixe de son solde de repo. Si
construire l'accès coûte $F_b > F_r$ mais $u_b < u_r$ par rapport à la location, construire est
moins cher exactement lorsque le solde de repo dépasse
\[
B^* = \frac{F_b - F_r}{u_r - u_b}.
\]
\end{proposition}

\begin{proof}
Les coûts annuels valent $F_r + u_r B$ et $F_b + u_b B$~; le second est plus petit lorsque $(u_r -
u_b)B > F_b - F_r$.
\end{proof}

\begin{example}[Trois voies]\label{ex:m2:getting-access-rates-credit:routes}
Dans le modèle illustratif du chapitre, la location coûte 0.5~million de dollars par an de frais fixes et 11.5
points de base par an par dollar de solde de repo, y compris les swaps et la marge financée~;
l'achat coûte 8~millions de dollars de frais fixes et 5.9 points de base, y compris le capital
dans le fonds de compensation~; le bilatéral coûte 0.2 million et 18.5 points de base. Un fonds disposant de 5
milliards de dollars de repo et de 10 milliards de swaps paie 6.25~millions de dollars par an en location, 10.95~millions
à l'achat et 9.45~millions en bilatéral. Le bilatéral est moins cher que la location seulement en dessous de 0.43
milliard de dollars de repo~; l'achat est moins cher que la location au-delà de 13.4 milliards
de dollars (\cref{fig:m2:getting-access-rates-credit:costs}).
\end{example}

\begin{omfigure}
\begin{tikzpicture}
\begin{axis}[width=11cm,height=5cm,
  xlabel={solde de repo (Md USD), swaps deux fois plus importants},ylabel={coût annuel (M USD)},
  tick label style={font=\small},label style={font=\small},
  xmin=0,xmax=30,ymin=0,ymax=40,grid=major,grid style={gray!25},
  legend columns=3,legend style={font=\footnotesize,at={(0.5,-0.3)},anchor=north,draw=none,fill=none,/tikz/every even column/.append style={column sep=8pt}},
  table/col sep=comma]
\addplot[omDef,very thick] table[x=repo,y=rent]{figdata/markets-2/28-getting-access-rates-credit/costs.csv};
\addplot[omThm,very thick] table[x=repo,y=buy]{figdata/markets-2/28-getting-access-rates-credit/costs.csv};
\addplot[omProp,very thick,dashed] table[x=repo,y=bilateral]{figdata/markets-2/28-getting-access-rates-credit/costs.csv};
\addplot[only marks,mark=*,mark size=1.8pt,black,forget plot] coordinates {(13.39,15.90)};
\node[font=\footnotesize,anchor=north west] at (axis cs:13.8,15.2) {seuil de rentabilité 13.4};
\legend{{location (sponsorisée, courtier)},achat (direct),bilatéral}
\end{axis}
\end{tikzpicture}

\omcaption{Coût annuel d'accès au repo de Trésor et aux swaps en fonction du solde de repo,
avec un notionnel de swaps deux fois plus élevé, pour trois voies~: location (\omterm{def:m2:repo-and-specials:sponsored}{repo sponsorisé} et
\omterm{def:m2:getting-access-rates-credit:cb}{courtier compensateur} de swaps), achat (adhésions directes) et bilatéral. L'achat est plus coûteux au-delà de 13.4
milliards de dollars de repo. Coûts illustratifs~; données~: tutoriel du chapitre.}\label{fig:m2:getting-access-rates-credit:costs}
\end{omfigure}

Le seuil de rentabilité ne vaut que par ses données d'entrée, et le plus incertain est le
spread que facture le parrain, qui dépend de son bilan et de son appétit
(\cref{fig:m2:getting-access-rates-credit:sens}).

\begin{omfigure}
\begin{tikzpicture}
\begin{axis}[width=11cm,height=4.6cm,
  xlabel={spread du parrain sur le solde de repo (pb par an)},ylabel={seuil de rentabilité (Md USD)},
  tick label style={font=\small},label style={font=\small},
  xmin=4,xmax=12,ymin=0,ymax=50,grid=major,grid style={gray!25},
  table/col sep=comma]
\addplot[omDef,very thick] table[x=spread,y=breakeven]{figdata/markets-2/28-getting-access-rates-credit/sensitivity.csv};
\addplot[only marks,mark=*,mark size=1.8pt,black] coordinates {(5,28.85) (8,13.39) (12,7.81)};
\end{axis}
\end{tikzpicture}

\omcaption{Le solde de repo au-delà duquel l'adhésion directe coûte moins cher que le parrainage,
en fonction du spread du parrain, les autres coûts étant inchangés. À 5 points de base, il est de 28.8
milliards de USD, à 8 points de base 13.4 et à 12 points de base 7.8. Illustratif~; données~: le
tutoriel du chapitre.}\label{fig:m2:getting-access-rates-credit:sens}
\end{omfigure}

\section{Tutoriel~: trois façons d'accéder au repo et aux swaps}\label{tut:m2:getting-access-rates-credit:cost}

\begin{tutorial}
\textbf{Objectif.} Chiffrer trois routes d'accès au repo sur \omterm{def:m2:money-markets:bill}{bons du Trésor} et aux swaps pour un fonds de taille moyenne, trouver les
tailles de seuil de rentabilité et voir comment elles évoluent avec le spread du parrain et le coût fixe de
l'adhésion.
\textbf{État final~:} \cref{fig:m2:getting-access-rates-credit:costs,fig:m2:getting-access-rates-credit:sens},
\cref{ex:m2:getting-access-rates-credit:routes} et les chiffres du problème du week-end.
\begin{tutsteps}
\item \textbf{Le modèle}~: routes, coûts annuels, coûts unitaires, seuils de rentabilité.
\omcode{code/firm/clearcost/firm_clearcost.py}{12}{51}{Le modèle de coût d'accès.}\label{lst:m2:getting-access-rates-credit:model}
\item \textbf{Exécuter} \texttt{access\_cost\_demo.tutorial()},
\texttt{access\_cost\_demo.problem()} et \texttt{fig\_access\_cost.py}.
\end{tutsteps}
\textbf{Pour aller plus loin.} Faire croître le spread du parrain avec le solde du fonds, comme le ferait
le coût de bilan d'une banque, et vérifier si le seuil de rentabilité existe encore~; puis ajouter
la valeur de la compensation (netting) dont bénéficie un membre direct sur l'ensemble de ses contreparties.
\end{tutorial}

\section{Construction~: le modèle de coût d'accès}\label{bld:m2:getting-access-rates-credit:clearcost}

\begin{build}
\bfield{Rôle} Le mini-fonds décide de la manière dont il accède à chaque marché et
réexamine cette décision à mesure qu'il grandit. \bfield{Interface} \texttt{Route(name, fixed, repo\_bp, swap\_bp, repo\_margin,
swap\_margin, fund\_share, capital\_rate)}~; \texttt{annual\_cost(route, repo, swaps,
funding)}~; \texttt{variable\_cost}~; \texttt{unit\_cost(route, swaps\_per\_repo,
funding)}~; \texttt{breakeven(a, b, swaps\_per\_repo, funding)}~; \texttt{cheapest(routes,
repo, swaps, funding)}.
\bfield{Règles} Coûts annuels en dollars~; frais en points de base du solde ou du notionnel~;
marge financée à un spread~; capital du fonds de défaillance à un coût du capital~; coûts linéaires
en fonction de la taille du livre au-delà de la partie fixe.
\bfield{Tests d'acceptation} \texttt{code/firm/clearcost/tests/}~: le coût est fixe plus
linéaire~; au seuil de rentabilité, les deux routes coûtent le même prix, en dessous la location est
moins chère et au-dessus l'achat l'est~; une route dominée n'a pas de seuil de rentabilité.
\bfield{Pour aller plus loin} Coûts croissant avec la taille~; valeur de la compensation et de l'anonymat~;
expositions d'un membre direct au fonds de défaillance et à la répartition des pertes~; coûts de stress lorsqu'un
parrain se retire (livre~16, sur la société).
\end{build}

\begin{omsources}
\item ISDA, ``Legal guidelines for smart derivatives contracts~: the ISDA Master
Agreement'', février 2019.
\item ICMA, Global Master Repurchase Agreement~; GLEIF, introducing the legal entity
identifier.
\item SEC, Release No.~34-101694 (21 novembre 2024) approuvant les modifications des règles d'accès de FICC~;
DTCC, FICC Sponsored Service.
\item 17 CFR 37.202, exigences d'accès pour les swap execution facilities (SEF).
\end{omsources}

\section{Exercices}

\begin{exercise}[$\star$]\label{exo:m2:getting-access-rates-credit:1}
Que fait la compensation à la clôture (close-out netting), et pourquoi un dealer se soucie-t-il
de son caractère exécutoire dans le pays de son client~?
\end{exercise}

\begin{exercise}[$\star$]\label{exo:m2:getting-access-rates-credit:2}
Énumérer les documents dont un fonds a besoin avant son premier swap compensé.
\end{exercise}

\begin{exercise}[$\star$]\label{exo:m2:getting-access-rates-credit:3}
À quoi sert un LEI, et à quelles deux questions ses données de référence répondent-elles~?
\end{exercise}

\begin{exercise}[$\star\star$]\label{exo:m2:getting-access-rates-credit:4}
La location coûte 0.5 million par an et 11.5 points de base~; l'achat coûte 8 millions
et 5.9 points de base. Calculer le seuil de rentabilité à partir de la formule.
\end{exercise}

\begin{exercise}[$\star\star$]\label{exo:m2:getting-access-rates-credit:5}
Quelle est la différence entre la compensation sponsorisée et un service de compensation
par agent (done-away), du point de vue du fonds~?
\end{exercise}

\begin{exercise}[$\star\star$]\label{exo:m2:getting-access-rates-credit:6}
Pourquoi l'\omterm{def:m2:swap-clearing-and-the-ccp-basis:mandate}{obligation de compensation} sur \omterm{def:m2:money-markets:bill}{bons du Trésor} modifie-t-elle la comparaison
entre le \omterm{def:m2:repo-and-specials:triparty}{repo bilatéral} et les routes compensées~?
\end{exercise}

\begin{exercise}[$\star\star\star$]\label{exo:m2:getting-access-rates-credit:7}
\emph{Programmation.} Supprimer les swaps (un fonds ne faisant que du repo) et trouver le
seuil de rentabilité entre location et achat. Puis porter le coût fixe de l'adhésion à 12 millions de USD, les swaps
étant réintroduits. Expliquer les deux changements.
\end{exercise}

\begin{exercise}[$\star\star\star$]\label{exo:m2:getting-access-rates-credit:8}
\emph{Repérer l'erreur.} « L'adhésion directe est toujours moins chère à long terme~: les frais
sont plus faibles, et les coûts fixes ne sont payés qu'une seule fois. » La corriger.
\end{exercise}

\section{Problème~: construire ou louer}

\begin{problem}[{Problème du week-end --- quand l'adhésion directe est rentable}]\label{pb:m2:getting-access-rates-credit:1}
Le livre de taux d'un fonds dispose d'un solde de repo $B$ et d'un notionnel de swaps $2B$. Louer l'accès
(\omterm{def:m2:repo-and-specials:sponsored}{repo sponsorisé}, \omterm{def:m2:getting-access-rates-credit:cb}{courtier compensateur} de swaps) coûte 0.5 million de USD par an, 8 points de base
par an sur le solde de repo, 0.5 sur le notionnel de swaps, avec une marge de 2~\% du repo et de 1.5~\%
des swaps, financée à 0.5~\% par an. L'acheter (adhésion directe) coûte 8 millions de USD
par an, 1 et 0.2 points de base, la même marge, et une contribution au fonds de défaillance de 0.2~\%
du solde de repo, qui coûte 10~\% par an en capital. La voie bilatérale, non compensée, coûte 0.2
million, 12 et 1 points de base, avec une marge de 3~\% sur les deux. Tous les chiffres sont illustratifs.

\medskip\noindent\textbf{Partie I --- Coûts unitaires.}
\begin{enumerate}
\item Calculer le coût variable de chaque route par dollar de solde de repo, en points de base.
\item Que coûte chaque route pour 5 milliards de USD de repo et 10 milliards de swaps~?
\item Laquelle est la moins chère à cette taille~?
\item En dessous de quelle taille la voie bilatérale est-elle la moins chère~?
\item Pourquoi la contribution au fonds de défaillance compte-t-elle comme un coût~?
\end{enumerate}

\medskip\noindent\textbf{Partie II --- Le seuil de rentabilité.}
\begin{enumerate}[resume]
\item Trouver le solde de repo au-delà duquel acheter est plus avantageux que louer.
\item Que coûte chaque route à cette taille~?
\item Qu'économiserait l'achat pour 20~milliards de USD de repo~?
\item Comment le seuil de rentabilité change-t-il pour un fonds qui ne fait que du repo~?
\item Comment change-t-il si l'adhésion coûte 12~millions par an~?
\end{enumerate}

\medskip\noindent\textbf{Partie III --- Sensibilité.}
\begin{enumerate}[resume]
\item Quel est le seuil de rentabilité si le sponsor facture 5~points de base~? 12~?
\item Pourquoi la fourchette du sponsor est-elle l'hypothèse la plus incertaine~?
\item Quels coûts de l'adhésion directe le modèle omet-il~?
\item Que coûte la location que le modèle ne montre pas~?
\item Comment l'\omterm{def:m2:swap-clearing-and-the-ccp-basis:mandate}{obligation de compensation} modifie-t-elle la voie bilatérale~?
\end{enumerate}

\medskip\noindent\textbf{Partie IV --- Jugement.}
\begin{enumerate}[resume]
\item Le fonds prévoit de passer de 5 à 20 milliards en trois ans. Quand doit-il
commencer à se préparer~?
\item Pourquoi un fonds pourrait-il louer même au-dessus du seuil de rentabilité~?
\item Quels documents doivent être en place quelle que soit la voie choisie~?
\item Énoncer le \emph{résultat nommé}~: la taille du bilan à partir de laquelle l'adhésion directe
l'emporte sur un \omterm{def:m2:getting-access-rates-credit:cb}{courtier compensateur}, et son coût.
\item En une phrase~: qu'achète un fonds lorsqu'il devient membre~?
\end{enumerate}
\end{problem}

\section{Questions d'entretien}

\begin{interviewq}[$\star$ \iqroles{trader, développeur}]\label{iq:m2:getting-access-rates-credit:1}
Qu'est-ce qu'un \omterm{def:m2:getting-access-rates-credit:isda}{contrat-cadre ISDA}, et pourquoi la compensation de close-out est-elle importante~?
\end{interviewq}

\begin{interviewq}[$\star$ \iqroles{trader, banque}]\label{iq:m2:getting-access-rates-credit:2}
Comment un hedge fund accède-t-il au repo sur \omterm{def:m2:money-markets:bill}{bons du Trésor} compensé~?
\end{interviewq}

\begin{interviewq}[$\star\star$ \iqroles{chercheur, risque}]\label{iq:m2:getting-access-rates-credit:3}
Construisez un modèle pour choisir entre un \omterm{def:m2:getting-access-rates-credit:cb}{courtier compensateur} et l'adhésion directe.
\end{interviewq}

\begin{interviewq}[$\star\star$ \iqroles{trader}]\label{iq:m2:getting-access-rates-credit:4}
Que fait un \omterm{def:m2:getting-access-rates-credit:cb}{courtier compensateur} en swaps, et que peut-il vous faire en période de stress~?
\end{interviewq}

\begin{interviewq}[$\star\star$ \iqroles{trader, chercheur}]\label{iq:m2:getting-access-rates-credit:5}
Comment les non-banques accèdent-elles au marché interdealer des \omterm{def:m2:money-markets:bill}{bons du Trésor}~?
\end{interviewq}

\begin{interviewq}[$\star\star\star$ \iqroles{développeur}]\label{iq:m2:getting-access-rates-credit:6}
Concevez le système d'intégration qui suit, pour chaque contrepartie, les documents signés,
leurs termes clés et ce qu'ils permettent à la société de négocier.
\end{interviewq}
